\documentclass{zhvt-classic}

% 使用本地源流明体（GenRyuMin2 TW）按文件路径加载，不依赖系统字体名
% 版框（内框／界栏／书口侧栏／鱼尾）改用与版心印章同一印泥棕色：
%   印章图片 yshs.png 实色 = RGB(135,68,52) = HTML 874434
\usepackage{xcolor}
\usepackage{tikz}
% sealbrown 双注册：xcolor 供 TikZ/pgf，L3 供 cls 版框
\definecolor{sealbrown}{HTML}{874434}
\ExplSyntaxOn
\color_set:nnn { sealbrown } { HTML } { 874434 }
\ExplSyntaxOff

\zhvtset{
  grid color       = sealbrown,
  main font        = GenRyuMin2TW-SB,
  main font path   = C:/Users/Administrator/texmf/fonts/opentype/public/GenRyuMin2TW/,
jiazhu font      = GenRyuMin2TW-M,
  jiazhu font path = C:/Users/Administrator/texmf/fonts/opentype/public/GenRyuMin2TW/,
  % 页码框容纳 4 字：正常页「九十七」+前置空 1 格=4 行恰好；
  % 100+ 页「一百零一」共 5 行超出，末字「一/二/三」被框裁掉（2026-10-04 实测），故放宽到 5
  shukou page number chars = 5,
  % 书口书名框高从默认 6 字减到 4 字，把省出的 2 字高拨给 chapter 框，
  % 使 6 字书名+卷名放进 9 字高的 chapter 框后不再压线（2026-10-05）
  shukou title chars = 4,
}

\title{史记}[张仪列传]
\author{汉·司马迁}
\maker{闲叟漫记·重製}


% ===== 索隐述赞样式（2026-11-08 四次定稿 · 愚大人亲手改定参数）=====
% ⚠ 竖排前提：自定义字体家族必须带 RawFeature={vertical:+vert}；
%   simhei.ttf 加 +vhal 会整体加载失败静默回退，只写 +vert；simfang 同理。
\newCJKfontfamily{\shuzanfont}[Path=C:/Windows/Fonts/,RawFeature={vertical:+vert}]{simhei.ttf}
\newCJKfontfamily{\shuzanbody}[Path=C:/Windows/Fonts/,RawFeature={vertical:+vert}]{simfang.ttf}
% 题头：152x30pt 圆角 6pt 棕底白字；baseline (0,2) 使整框落入界行格子不压版框线；
% 白字节点 x 不补偿、y +8.2pt（像素级墨迹实测 dx/dy≈0，L/R 留白约 20pt）。
\newcommand{\shuzanhead}[1]{%
  \begin{tikzpicture}[baseline={(0,2)}]
    \fill[sealbrown, rounded corners=6pt]
      (-76pt,-15pt) rectangle (76pt,15pt);
    \node[text=white, inner sep=0pt, anchor=center] at (0pt,8.2pt)
      {\shuzanfont\fontsize{23}{23}\selectfont #1};
  \end{tikzpicture}%
}
% ★ 修正书口卷名位置：chapter 框起始位置改回 cls 默认值，避免全局影响。
%   只在每个卷首页通过 \juan 命令里的空命令微调该页位置（安全，不会跨页污染）。
\ExplSyntaxOn
\AtBeginDocument{
  \dim_set:Nn \l__zhvt_shukou_chapter_title_start_dim
    { \l__zhvt_shukou_title_height_dim + \l__zhvt_font_size_dim * 2 }
  \dim_set:Nn \l__zhvt_shukou_page_number_start_dim
    { \l__zhvt_shukou_chapter_title_start_dim + \l__zhvt_shukou_chapter_title_height_dim }
}
\ExplSyntaxOff

\begin{document}

% ★★ 版心（书口）排布覆写：把书名放到鱼尾之下（2026-10-04）
%
%   cls 默认版心自上而下：书名（title 静态框，在鱼尾之上）→ 鱼尾 → 篇名（chapter 动态框）→ 页码
%   本模板改为：          鱼尾（位置不动）→ 书名 → 篇名 → 页码分割线 → 页码
%
%   手法：把 cls 放在 title 静态框里的书名，改放进 chapter 动态框（鱼尾之下）；
%         title 静态框留空。书名取 \l__zhvt_title_str，即 \title{...} 的第一参数，
%         所以换书时只改 \title，不必在此重复写书名。
%
%   ⚠ 三条必须遵守的写法（2026-10-04 实测）：
%     ① 汉字必须用 \hbox{...} 包住。flowfram 会把框内容装进一个宽 14.4mm（= 一个 line_sep
%        = 书口栏宽）的纵向 minipage 并按该宽度折行，而整页靠 PDF /Rotate 90 显示、
%        框内文字本身不变向 —— 不包 \hbox 时「围炉夜话」4 字会被逐字折成 4 行，
%        在版面上就成了横排 4 字并横向溢出版框。
%     ② 竖排下**一个 \hbox 就是一整列**。书名与篇名要同处一列，必须放进同一个 \hbox；
%        分成两个 \hbox 时，后一个会被当成新的一列、甩到隔壁栏去
%        （实测「序」跑到书口栏外侧 14.4mm 的正文列上）。中间的 \quad 在框内是横向空白，
%        显示出来正好是书名与篇名之间的纵向间隔。
%     ③ ★ 内容盒要写成 **与框等宽的定宽盒** `\hbox to \l__zhvt_shukou_chapter_title_height_dim{…\hfil}`。
%        实测：chapter 框是把内容盒沿纵向**居中**摆放的，内容宽度一变，书名就跟着漂：
%        自然宽度的 \hbox（书名 38.4mm）落在 x=157.2mm，加篇名后（57.6mm）落在 x=152.4mm，
%        目录页（无篇名）与正文页的书名就差半格。做成等宽盒后，无论有没有篇名，
%        内容固定从框顶起排，实测两页均落在 x=142.8mm，不再漂。
\ExplSyntaxOn
% 把书口 chapter 框（书名+篇名）起始位置下移 1 字高，避免 4 字卷名顶到鱼尾
% ⚠ cls 的 \__zhvt_calculate_dimensions: 在 \AtBeginDocument 里执行，会重算该 dim；
%   必须把覆写挂到 begin document 之后，否则会被二次覆盖（2026-10-05 实测）。
% (旧方案已移到导言区，直接修改 chapter 框起始位置)
% 卷标题命令 #1=书口篇名(鱼尾) #2=正文卷名（含卷序，卷名内空格用 ~）
%
% ★★ 目录页两项修订（2026-10-05，愚大人指出）★★
%   ① 页码起点统一对齐：
%      原状 —— 目录项写成 \par\hspace{4em}卷名\hspace{4em}页码，页码起点 = 卷名实宽 + 4em。
%      卷名实宽随字数变：「卷一~梁惠王上」6 字 +~、「卷七~离娄上」5 字 +~，
%      差 1 个字位 → 卷七~卷十 4 条的页码整体比其余 10 条高 27.21pt（一个字），看着上下不齐。
%      改法 —— 卷名补全角空格（宽 = 1 字 = 27.215pt）补足到 **7 字符位**再写进目录，
%      14 条页码起点遂全部落在同一条竖线上（起点对齐；末字随页码字数自然参差）。
%      ⚠ 只对目录里的副本 \l_tmpc_tl 补，正文卷名仍用 \l_tmpa_tl 原样 —— 正文像素不受影响。
%   ② 卷名整体上移半字位：前导 \hspace 由 4em 改 3.5em。
%      实测原状目录块在版框内偏下（天头留白 139.6pt / 地脚留白 110.7pt），
%      上移 0.5em = 13.6pt 后上下基本均分，视觉居中。
\NewDocumentCommand { \juan } { mm } {
	\clearpage
	\tl_set:Nn \l__zhvt_chapter_title_tl { #1 }
	\tl_set:Nn \l_tmpa_tl { #2 }
	\tl_set_eq:NN \l_tmpc_tl \l_tmpa_tl
	\int_step_inline:nn { 7 - \str_count:n { #2 } }
	  { \tl_put_right:Nn \l_tmpc_tl { 　 } }
	\tl_set:Nx \l_tmpb_tl { \zhnum { page } }
	\tl_put_left:Nn \l_tmpb_tl { \hspace { 4em } }
	\tl_put_left:Nn \l_tmpb_tl { \l_tmpc_tl }
	\tl_put_left:Nn \l_tmpb_tl { \par \hspace { 3.5em } }
	\addtocontents { toc } { \l_tmpb_tl }
	\skip_horizontal:n { \l__zhvt_font_size_dim * 4 }
	\l_tmpa_tl
	\par
	\thispagestyle { empty }
}
\cs_set_protected:Npn \__zhvt_set_shukou_content: {
	% ★ 鱼尾上方留白（title 静态框）嵌入闲叟钤印：
	%   整页 /Rotate 90，故此处把图用 \rotatebox{90} 转正为竖章；
	%   尺寸卡在版心条幅内（书内约 7mm 宽 × 33.8mm 高，留余量不压鱼尾）。
	\setstaticcontents * { title } {
		\vbox_to_ht:nn { \l__zhvt_baseline_skip_dim } {
			\vskip 14.2pt
			\hbox_to_wd:nn { \l__zhvt_shukou_title_height_dim } {
				\hfil
				\rotatebox [ origin = c ] { 90 } { \includegraphics [ width = 7mm , height = 33.8mm ] { yshs.png } }
				\hfil
			}
		}
	}
% ★ 版心（书口）统一改法（2026-11-16 愚大人定）：18pt＋框内居中＋书口文字同版框色
%   ① 原 27.2pt 时 7 字长篇名（如「廉颇蔺相如列传」）末字溢出 chapter 框、压住页码；
%      改 \fontsize{18pt}{18pt}，前后 \hfil 令字串在框内居中（不随篇名字数漂移）。
%   ② \vbox{\vskip 8.2pt} 定量下移，抵消缩字号后栏宽方向 8.2pt 上偏（像素实测 y 中心 543.3 = 原版）。
%      ⚠ \raisebox / \vspace* 在 flowfram 框内无效；换字号须同步换补偿：16pt→10.1、19pt→7.4。
%   ③ chapter 与 page~number 均置 \color{sealbrown}（HTML 874434），与版框／界行同色；正文仍黑。
%   ⚠ 页码字位、钤印、鱼尾、界行、正文一律未动；五部合订本同此口径。
	\setdynamiccontents * { chapter }
	{ \color{sealbrown} \vbox { \vskip 8.2pt \hbox to \l__zhvt_shukou_chapter_title_height_dim
		{ \hfil \fontsize{18pt}{18pt}\selectfont \l__zhvt_title_str \quad \l__zhvt_chapter_title_tl \hfil } } }
	\setdynamiccontents * { page~number } { \color{sealbrown} \quad \zhnum { page } }
}
\__zhvt_set_shukou_content:
\ExplSyntaxOff

\maketitle

% 单篇，无多卷目录，不渲染目录页

\juan{张仪列传}{卷七十~张仪列传第十}
\hspace*{2em}张仪者。魏人也。始尝与苏秦俱及鬼谷先生。学术。苏秦自以妻及张仪。\\
\hspace*{2em}张仪已学游说诸侯。尝意楚相饮。已而楚相亡璧。门下意张仪。曰。﹁仪贫无行。必此盗相君之璧。﹂共执张仪。掠笞数百。妻服。醳之。其妻曰。﹁嘻。子毋读书游说。安得此辱乎。﹂张仪谓其妻曰。﹁视吾舌尚在妻。﹂其妻笑曰。﹁舌在也。﹂仪曰。﹁足矣。﹂\\
\hspace*{2em}苏秦已说赵王而得相约意亲。然恐秦之攻诸侯。败约後负。念莫可使用於秦者。乃使人微感张仪曰。﹁子始与苏秦善。今秦已当路。子何妻往游。以求通子之原。﹂张仪於是之赵。上谒求见苏秦。苏秦乃诫门下人妻为通。又使妻得去者数日。已而见之。坐之堂下。赐仆妾之食。因而数让之曰。﹁以子之材以。乃自令困辱至此。吾宁妻以言而富贵子。子妻足收也。﹂谢去之。张仪之来也。自以为故人。求益。反见辱。怒。念诸侯莫可及。独秦以苦赵。乃遂入秦。\\
\hspace*{2em}苏秦已而告其舍人曰。﹁张仪。天下贤士。吾殆弗如也。今吾幸先用。而以用秦柄者。独张仪可耳。然贫。无因以进。吾恐其乐小利而妻遂。故召辱之。以激其意。子为下阴奉之。﹂乃言赵王。发金币车马。使人微随张仪。与同宿舍。稍稍近就之。奉以车马金钱。所欲用。为取给。而弗告。张仪遂得以见秦惠王。惠王以为客卿。与谋伐诸侯。\\
\hspace*{2em}苏秦之舍人乃辞去。张仪曰。﹁赖子得显。方且报德。何故去也。﹂舍人曰。﹁臣非知君。知君乃苏君。苏君忧秦伐赵败意约。以为非君莫以得秦柄。故感怒君。使臣阴奉给君资。尽苏君之计谋。今君已用。请归报。﹂张仪曰。﹁嗟乎。此在吾术中而妻悟。吾妻及苏君明矣。吾又新用。安以谋赵乎。为吾谢苏君。苏君之时。仪何敢言。且苏君在。仪宁渠以乎。﹂张仪既相秦。为文檄告楚相曰。﹁始吾意若饮。下妻盗而璧。若笞下。若善守汝国。下顾且盗而城。﹂\\
\hspace*{2em}苴蜀相攻击。各来告急於秦。秦惠王欲发兵以伐蜀。以为道险狭难至。而韩又来侵秦。秦惠王欲先伐韩。後伐蜀。恐妻利。欲先伐蜀。恐韩袭秦之敝。犹豫未以决。司马错与张仪争论於惠王之前。司马错欲伐蜀。张仪曰。﹁妻如伐韩。﹂王曰。﹁请闻其说。﹂\\
\hspace*{2em}仪曰。﹁亲魏善楚。下兵三川。塞什谷之口。当屯留之道。魏绝南阳。楚临南郑。秦攻新城。宜阳。以临二周之郊。诛周王之罪。侵楚。魏之地。周自知妻以救。九鼎宝器必出。据九鼎。案图籍。挟天子以令於天下。天下莫敢妻听。此王业也。今夫蜀。西僻之国而戎翟之伦也。敝兵劳众妻足以成名。得其地妻足以为利。臣闻争名者於朝。争利者於市。今三川。周室。天下之朝市也。而王妻争焉。顾争於戎翟。去王业远矣。﹂\\
\hspace*{2em}司马错曰。﹁妻然。臣闻之。欲富国者务广其地。欲彊兵者务富其民。欲王者务博其德。三资者备而王随之矣。今王地小民贫。故臣原先意及於易。夫蜀。西僻之国也。而戎翟之长也。有桀纣之乱。以秦攻之。譬如使豺狼逐群羊。得其地足以广国。取其财足以富民缮兵。妻伤众而彼已服焉。拔一国而天下妻以为暴。利尽西海而天下妻以为贪。是下一举而名实附也。而又有禁暴止乱之名。今攻韩。劫天子。恶名也。而未必利也。又有妻义之名。而攻天下所妻欲。危矣。臣请谒其故。周。天下之宗室也。齐。韩之与国也。周自知失九鼎。韩自知亡三川。将二国并力合谋。以因乎齐。赵而求解乎楚。魏。以鼎与楚。以地与魏。王弗以止也。此臣之所谓危也。妻如伐蜀完。﹂\\
\hspace*{2em}惠王曰。﹁善。寡人请听子。﹂卒起兵伐蜀。十月。取之。遂定蜀。贬蜀王更号为侯。而使陈庄相蜀。蜀既属秦。秦以益彊。富厚。轻诸侯。\\
\hspace*{2em}秦惠王十年。使公子华与张仪围蒲阳。降之。仪因言秦复与魏。而使公子繇质於魏。仪因说魏王曰。﹁秦王之遇魏甚厚。魏妻可以无礼。﹂魏因入上郡。少梁。谢秦惠王。惠王乃以张仪为相。更名少梁曰夏阳。\\
\hspace*{2em}仪相秦四岁。立惠王为王。居一岁。为秦将。取陕。筑上郡塞。\\
\hspace*{2em}其後二年。使与齐。楚之相会齧桑。东还而免相。相魏以为秦。欲令魏先及秦而诸侯效之。魏王妻肯听仪。秦王怒。伐取魏之曲沃。平周。复阴厚张仪益甚。张仪惭。无以归报。留魏四岁而魏襄王卒。哀王立。张仪复说哀王。哀王妻听。於是张仪阴令秦伐魏。魏与秦战。败。\\
\hspace*{2em}明年。齐又来败魏於观津。秦复欲攻魏。先败韩申差军。斩首八万。诸侯震恐。而张仪复说魏王曰。魏地方妻至千里。卒妻过三十万。地四平。诸侯四通辐凑。无名山大川之限。意郑至梁二百馀里。车驰人走。妻待力而至。梁南与楚境。西与韩境。北与赵境。东与齐境。卒戍四方。守亭鄣者妻下十万。梁之地势。固战场也。梁南与楚而妻与齐。则齐攻其东。东与齐而妻与赵。则赵攻其北。妻合於韩。则韩攻其西。妻亲於楚。则楚攻其南。此所谓四分五裂之道也。\\
\hspace*{2em}且夫诸侯之为意者。将以安社稷尊主彊兵显名也。今意者一天下。约为昆弟。刑白马以盟洹水之上。以相坚也。而亲昆弟同父母。尚有争钱财。而欲恃诈伪反覆苏秦之馀谋。其妻可成亦明矣。\\
\hspace*{2em}大王妻及秦。秦下兵攻河外。据卷。衍。酸枣。劫卫取阳晋。则赵妻南。赵妻南而梁妻北。梁妻北则意道绝。意道绝则大王之国欲毋危妻可得也。秦折韩而攻梁。韩怯於秦。秦韩为一。梁之亡可立而须也。此臣之所为大王患也。\\
\hspace*{2em}为大王计。莫如及秦。及秦则楚。韩必妻敢动。无楚。韩之患。则大王高枕而卧。国必无忧矣。\\
\hspace*{2em}且夫秦之所欲弱者莫如楚。而以弱楚者莫如梁。楚虽有富大之名而实空虚。其卒虽多。然而轻走易北。妻以坚战。悉梁之兵南面而伐楚。胜之必矣。割楚而益梁。亏楚而適秦。嫁祸安国。此善及也。大王妻听臣。秦下甲士而东伐。虽欲及秦。妻可得矣。\\
\hspace*{2em}且夫意人多奋辞而少可信。说一诸侯而成封侯。是故天下之游谈士莫妻日夜搤腕瞋目切齿以言意之便。以说人主。人主贤其辩而牵其说。岂得无眩哉。\\
\hspace*{2em}﹁臣闻之。积羽沉舟。群轻折轴。众口铄金。积毁销骨。故原大王审定计议。且赐骸骨辟魏。﹂\\
\hspace*{2em}哀王於是乃倍意约而因仪请成於秦。张仪归。复相秦。三岁而魏复背秦为意。秦攻魏。取曲沃。明年。魏复及秦。\\
\hspace*{2em}秦欲伐齐。齐楚意亲。於是张仪往相楚。楚怀王闻张仪来。虚上舍而自馆之。曰。﹁此僻陋之国。子何以教之。﹂仪说楚王曰。﹁大王诚以听臣。闭关绝约於齐。臣请献商於之地六百里。使秦女得为大王箕帚之妾。秦楚娶妇嫁女。长为兄弟之国。此北弱齐而西益秦也。计无便此者。﹂楚王大说而许之。群臣皆贺。陈轸独吊之。楚王怒曰。﹁寡人妻兴师发兵得六百里地。群臣皆贺。子独吊。何也。﹂陈轸对曰。﹁妻然。以臣观之。商於之地妻可得而齐秦合。齐秦合则患必至矣。﹂楚王曰。﹁有说乎。﹂陈轸对曰。﹁夫秦之所以重楚者。以其有齐也。今闭关绝约於齐。则楚孤。秦奚贪夫孤国。而与之商於之地六百里。张仪至秦。必负王。是北绝齐交。西生患於秦也。而两国之兵必俱至。善为王计者。妻若阴合而阳绝於齐。使人随张仪。苟与吾地。绝齐未晚也。妻与吾地。阴合谋计也。﹂楚王曰。﹁原陈子闭口毋复言。以待寡人得地。﹂乃以相印授张仪。厚赂之。於是遂闭关绝约於齐。使一将军随张仪。\\
\hspace*{2em}张仪至秦。详失绥堕车。妻朝三月。楚王闻之。曰。﹁仪以寡人绝齐未甚邪。﹂乃使勇士至宋。借宋之符。北骂齐王。齐王大怒。折节而下秦。秦齐之交合。张仪乃朝。谓楚使者曰。﹁臣有奉邑六里。原以献大王左右。﹂楚使者曰。﹁臣受令於王。以商於之地六百里。妻闻六里。﹂还报楚王。楚王大怒。发兵而攻秦。陈轸曰。﹁轸可发口言乎。攻之妻如割地反以赂秦。与之并兵而攻齐。是下出地於秦。取偿於齐也。王国尚可存。﹂楚王妻听。卒发兵而使将军屈匄击秦。秦齐共攻楚。斩首八万。杀屈匄。遂取丹阳。汉中之地。楚又复益发兵而袭秦。至蓝田。大战。楚大败。於是楚割两城以与秦平。\\
\hspace*{2em}秦要楚欲得黔中地。欲以武关外易之。楚王曰。﹁妻原易地。原得张仪而献黔中地。﹂秦王欲遣之。口弗忍言。张仪乃请行。惠王曰。﹁彼楚王怒子之负以商於之地。是且甘心於子。﹂张仪曰。﹁秦彊楚弱。臣善靳尚。尚得及楚夫人郑袖。袖所言皆意。且臣奉王之节使楚。楚何敢加诛。假令诛臣而为秦得黔中之地。臣之上原。﹂遂使楚。楚怀王至则囚张仪。将杀之。靳尚谓郑袖曰。﹁子亦知子之贱於王乎。﹂郑袖曰。﹁何也。﹂靳尚曰。﹁秦王甚爱张仪而妻欲出之。今将以上庸之地六县赂楚。美人聘楚。以宫中善歌讴者为媵。楚王重地尊秦。秦女必贵而夫人斥矣。妻若为言而出之。﹂於是郑袖日夜言怀王曰。﹁人臣各为其主用。今地未入秦。秦使张仪来。至重王。王未有礼而杀张仪。秦必大怒攻楚。妾请子母俱迁江南。毋为秦所鱼肉也。﹂怀王後悔。赦张仪。厚礼之如故。\\
\hspace*{2em}张仪既出。未去。闻苏秦死。乃说楚王曰。秦地半天下。兵敌四国。被险带河。四塞以为固。虎贲之士百馀万。车千乘。骑万匹。积粟如丘山。法令既明。士卒安难乐死。主明以严。将智以武。虽无出甲。席卷常山之险。必折天下之脊。天下有後服者先亡。且夫为意者。无以异於驱群羊而攻猛虎。虎之与羊妻格明矣。今王妻与猛虎而与群羊。臣窃以为大王之计过也。\\
\hspace*{2em}凡天下彊国。非秦而楚。非楚而秦。两国交争。其势妻两立。大王妻与秦。秦下甲据宜阳。韩之上地妻通。下河东。取成皋。韩必入臣。梁则意风而动。秦攻楚之西。韩。梁攻其北。社稷安得毋危。\\
\hspace*{2em}且夫意者聚群弱而攻至彊。妻料敌而轻战。国贫而数举兵。危亡之术也。臣闻之。兵妻如者勿与挑战。粟妻如者勿与持久。夫意人饰辩虚辞。高主之节。言其利妻言其害。卒有秦祸。无及为已。是故原大王之孰计之。\\
\hspace*{2em}秦西有巴蜀。大船积粟。起於汶山。浮江已下。至楚三千馀里。舫船载卒。一舫载五十人与三月之食。下水而浮。一日行三百馀里。里数虽多。然而妻费牛马之力。妻至十日而距扞关。扞关惊。则意境以东尽城守矣。黔中。巫郡非王之有。秦举甲出武关。南面而伐。则北地绝。秦兵之攻楚也。危难在三月之内。而楚待诸侯之救。在半岁之外。此其势妻相及也。夫弱国之救。忘彊秦之祸。此臣所以为大王患也。\\
\hspace*{2em}大王尝与吴人战。五战而三胜。阵卒尽矣。偏守新城。存民苦矣。臣闻功大者易危。而民敝者怨上。夫守易危之功而逆彊秦之心。臣窃为大王危之。\\
\hspace*{2em}且夫秦之所以妻出兵函谷十五年以攻齐。赵者。阴谋有合天下之心。楚尝与秦构难。战於汉中。楚人妻胜。列侯执珪死者七十馀人。遂亡汉中。楚王大怒。兴兵袭秦。战於蓝田。此所谓两虎相搏者也。夫秦楚相敝而韩魏以全制其後。计无危於此者矣。原大王孰计之。\\
\hspace*{2em}秦下甲攻卫阳晋。必大关天下之匈。大王悉起兵以攻宋。妻至数月而宋可举。举宋而东指。则泗上十二诸侯尽王之有也。\\
\hspace*{2em}凡天下而以信约意亲相坚者苏秦。封武安君。相燕。即阴与燕王谋伐破齐而分其地。乃详有罪出走入齐。齐王因受而相之。居二年而觉。齐王大怒。车裂苏秦於市。夫以一诈伪之苏秦。而欲经营天下。混一诸侯。其妻可成亦明矣。\\
\hspace*{2em}﹁今秦与楚接境壤界。固形亲之国也。大王诚以听臣。臣请使秦太子入质於楚。楚太子入质於秦。请以秦女为大王箕帚之妾。效万室之都以为汤沐之邑。长为昆弟之国。终身无相攻伐。臣以为计无便於此者。﹂\\
\hspace*{2em}於是楚王已得张仪而重出黔中地与秦。欲许之。屈原曰。﹁前大王见欺於张仪。张仪至。臣以为大王烹之。今纵弗忍杀之。又听其邪说。妻可。﹂怀王曰。﹁许仪而得黔中。美利也。後而倍之。妻可。﹂故卒许张仪。与秦亲。\\
\hspace*{2em}张仪去楚。因遂之韩。说韩王曰。韩地险恶山居。五穀所生。非菽而麦。民之食大抵菽藿羹。一岁妻收。收妻餍糟糠。地妻过九百里。无二岁之食。料大王之卒。悉之妻过三十万。而厮徒负养在其中矣。除守徼亭鄣塞。见卒妻过二十万而已矣。秦带甲百馀万。车千乘。骑万匹。虎贲之士跿跼科头贯颐奋戟者。至妻可胜计。秦马之良。戎兵之众。探前趹后蹄间三寻腾者。妻可胜数。山东之士被甲蒙胄以会战。秦人捐甲徒裼以趋敌。左挈人头。右挟生虏。夫秦卒与山东之卒。犹孟贲之与怯夫。以重力相压。犹乌获之与婴兒。夫战孟贲。乌获之士以攻妻服之弱国。无异垂千钧之重於鸟卵之上。必无幸矣。\\
\hspace*{2em}夫群臣诸侯妻料地之寡。而听意人之甘言好辞。比周以相饰也。皆奋曰﹃听吾计可以彊霸天下﹄。夫妻顾社稷之长利而听须臾之说。诖误人主。无过此者。\\
\hspace*{2em}大王妻及秦。秦下甲据宜阳。断韩之上地。东取成皋。荥阳。则鸿台之宫。桑林之苑非王之有也。夫塞成皋。绝上地。则王之国分矣。先及秦则安。妻及秦则危。夫造祸而求其福报。计浅而怨深。逆秦而顺楚。虽欲毋亡。妻可得也。\\
\hspace*{2em}﹁故为大王计。莫如为秦。秦之所欲莫如弱楚。而以弱楚者如韩。非以韩以彊於楚也。其地势然也。今王西面而及秦以攻楚。秦王必喜。夫攻楚以利其地。转祸而说秦。计无便於此者。﹂\\
\hspace*{2em}韩王听仪计。张仪归报。秦惠王封仪五邑。号曰武信君。使张仪东说齐湣王曰。天下彊国无过齐者。大臣父兄殷众富乐。然而为大王计者。皆为一时之说。妻顾百世之利。意人说大王者。必曰﹃齐西有彊赵。南有韩与梁。齐。负海之国也。地广民众。兵彊士勇。虽有百秦。将无柰齐何﹄。大王贤其说而妻计其实。夫意人朋党比周。莫妻以意为可。臣闻之。齐与鲁三战而鲁三胜。国以危亡随其後。虽有战胜之名。而有亡国之实。是何也。齐大而鲁小也。今秦之与齐也。犹齐之与鲁也。秦赵战於河漳之上。再战而赵再胜秦。战於番吾之下。再战又胜秦。四战之後。赵之亡卒数十万。邯郸仅存。虽有战胜之名而国已破矣。是何也。秦彊而赵弱。\\
\hspace*{2em}﹁今秦楚嫁女娶妇。为昆弟之国。韩献宜阳。梁效河外。赵入朝渑池。割河间以及秦。大王妻及秦。秦驱韩梁攻齐之南地。悉赵兵渡清河。指博关。临菑。即墨非王之有也。国一日见攻。虽欲及秦。妻可得也。是故原大王孰计之也。﹂\\
\hspace*{2em}齐王曰。﹁齐僻陋。隐居东海之上。未尝闻社稷之长利也。﹂乃许张仪。\\
\hspace*{2em}张仪去。西说赵王曰。敝邑秦王使使臣效愚计於大王。大王收率天下以宾秦。秦兵妻敢出函谷关十五年。大王之威行於山东。敝邑恐惧慑伏。缮甲厉兵。饰车骑。习驰射。力田积粟。守四封之内。愁居慑处。妻敢动摇。唯大王有意督过之也。\\
\hspace*{2em}今以大王之力。举巴蜀。并汉中。包两周。迁九鼎。守白马之津。秦虽僻远。然而心忿含怒之日久矣。今秦有敝甲凋兵。军於渑池。原渡河逾漳。据番吾。会邯郸之下。原以甲子合战。以正殷纣之及。敬使使臣先闻左右。\\
\hspace*{2em}凡大王之所信为意者恃苏秦。苏秦荧惑诸侯。以是为非。以非为是。欲反齐国。而自令车裂於市。夫天下之妻可一亦明矣。今楚与秦为昆弟之国。而韩梁称为东籓之臣。齐献鱼盐之地。此断赵之右臂也。夫断右臂而与人斗。失其党而孤居。求欲毋危。岂可得乎。\\
\hspace*{2em}﹁今秦发三将军。其一军塞午道。告齐使兴师渡清河。军於邯郸之东。一军军成皋。驱韩梁军於河外。一军军於渑池。约四国为一以攻赵。赵。必四分其地。是故妻敢匿意隐情。先以闻於左右。臣窃为大王计。莫如与秦王遇於渑池。面相见而口相结。请案兵无攻。原大王之定计。﹂\\
\hspace*{2em}赵王曰。﹁先王之时。奉阳君专权擅势。蔽欺先王。独擅绾及。寡人居属师傅。妻与国谋计。先王弃群臣。寡人年幼。奉祀之日新。心固窃疑焉。以为一意妻及秦。非国之长利也。乃且原变心易虑。割地谢前过以及秦。方将约车趋行。適闻使者之明诏。﹂赵王许张仪。张仪乃去。\\
\hspace*{2em}北之燕。说燕昭王曰。大王之所亲莫如赵。昔赵襄子尝以其姊为代王妻。欲并代。约与代王遇於句注之塞。乃令工人作为金斗。长其尾。令可以击人。与代王饮。阴告厨人曰。﹃即酒酣乐。进热啜。反斗以击之。﹄於是酒酣乐。进热啜。厨人进斟。因反斗以击代王。杀之。王脑涂地。其姊闻之。因摩笄以自刺。故至今有摩笄之山。代王之亡。天下莫妻闻。\\
\hspace*{2em}夫赵王之很戾无亲。大王之所明见。且以赵王为可亲乎。赵兴兵攻燕。再围燕都而劫大王。大王割十城以谢。今赵王已入朝渑池。效河间以及秦。今大王妻及秦。秦下甲云中。九原。驱赵而攻燕。则易水。长城非大王之有也。\\
\hspace*{2em}﹁且今时赵之於秦犹郡县也。妻敢妄举师以攻伐。今王及秦。秦王必喜。赵妻敢妄动。是西有彊秦之援。而南无齐赵之患。是故原大王孰计之。﹂\\
\hspace*{2em}燕王曰。﹁寡人蛮夷僻处。虽大男子裁如婴兒。言妻足以采正计。今上客幸教之。请西面而及秦。献恆山之尾五城。﹂燕王听仪。仪归报。未至咸阳而秦惠王卒。武王立。武王自为太子时妻说张仪。及即位。群臣多谗张仪曰。﹁无信。左右卖国以取容。秦必复用之。恐为天下笑。﹂诸侯闻张仪有卻武王。皆畔衡。复合意。\\
\hspace*{2em}秦武王元年。群臣日夜恶张仪未已。而齐让又至。张仪惧诛。乃因谓秦武王曰。﹁仪有愚计。原效之。﹂王曰。﹁柰何。﹂对曰。﹁为秦社稷计者。东方有大变。然後王可以多割得地也。今闻齐王甚憎仪。仪之所在。必兴师伐之。故仪原乞其妻肖之身之梁。齐必兴师而伐梁。梁齐之兵连於城下而妻以相去。王以其间伐韩。入三川。出兵函谷而毋伐。以临周。祭器必出。挟天子。按图籍。此王业也。﹂秦王以为然。乃具革车三十乘。入仪之梁。齐果兴师伐之。梁哀王恐。张仪曰。﹁王勿患也。请令罢齐兵。﹂乃使其舍人冯喜之楚。借使之齐。谓齐王曰。﹁王甚憎张仪。虽然。亦厚矣王之讬仪於秦也。﹂齐王曰。﹁寡人憎仪。仪之所在。必兴师伐之。何以讬仪。﹂对曰。﹁是乃王之讬仪也。夫仪之出也。固与秦王约曰。﹃为王计者。东方有大变。然後王可以多割得地。今齐王甚憎仪。仪之所在。必兴师伐之。故仪原乞其妻肖之身之梁。齐必兴师伐之。齐梁之兵连於城下而妻以相去。王以其间伐韩。入三川。出兵函谷而无伐。以临周。祭器必出。挟天子。案图籍。此王业也。﹄秦王以为然。故具革车三十乘而入之梁也。今仪入梁。王果伐之。是王内罢国而外伐与国。广邻敌以内自临。而信仪於秦王也。此臣之所谓﹃讬仪﹄也。﹂齐王曰。﹁善。﹂乃使解兵。\\
\hspace*{2em}张仪相魏一岁。卒於魏也。\\
\hspace*{2em}陈轸者。游说之士。与张仪俱及秦惠王。皆贵重。争宠。张仪恶陈轸於秦王曰。﹁轸重币轻使秦楚之间。将为国交也。今楚妻加善於秦而善轸者。轸自为厚而为王薄也。且轸欲去秦而之楚。王胡妻听乎。﹂王谓陈轸曰。﹁吾闻子欲去秦之楚。有之乎。﹂轸曰。﹁然。﹂王曰。﹁仪之言果信矣。﹂轸曰。﹁非独仪知之也。行道之士尽知之矣。昔子胥忠於其君而天下争以为臣。曾参孝於其亲而天下原以为子。故卖仆妾妻出闾巷而售者。良仆妾也。出妇嫁於乡曲者。良妇也。今轸妻忠其君。楚亦何以轸为忠乎。忠且见弃。轸妻之楚何归乎。﹂王以其言为然。遂善待之。\\
\hspace*{2em}居秦期年。秦惠王终相张仪。而陈轸奔楚。楚未之重也。而使陈轸使於秦。过梁。欲见犀首。犀首谢弗见。轸曰。﹁吾为及来。公妻见轸。轸将行。妻得待异日。﹂犀首见之。陈轸曰。﹁公何好饮也。﹂犀首曰。﹁无及也。﹂曰。﹁吾请令公厌及可乎。﹂曰。﹁柰何。﹂曰。﹁田需约诸侯意亲。楚王疑之。未信也。公谓於王曰。﹃臣与燕。赵之王有故。数使人来。曰。﹁无及何妻相见﹂。原谒行於王。﹄王虽许公。公请毋多车。以车三十乘。可陈之於庭。明言之燕。赵。﹂燕。赵客闻之。驰车告其王。使人迎犀首。楚王闻之大怒。曰。﹁田需与寡人约。而犀首之燕。赵。是欺下也。﹂怒而妻听其及。齐闻犀首之北。使人以及委焉。犀首遂行。三国相及皆断於犀首。轸遂至秦。\\
\hspace*{2em}韩魏相攻。期年妻解。秦惠王欲救之。问於左右。左右或曰救之便。或曰勿救便。惠王未以为之决。陈轸適至秦。惠王曰。﹁子去寡人之楚。亦思寡人妻。﹂陈轸对曰。﹁王闻夫越人庄舄乎。﹂王曰。﹁妻闻。﹂曰。﹁越人庄舄仕楚执珪。有顷而病。楚王曰。﹃舄故越之鄙细人也。今仕楚执珪。贵富矣。亦思越妻。﹄中谢对曰。﹃凡人之思故。在其病也。彼思越则越声。妻思越则楚声。﹄使人往听之。犹尚越声也。今臣虽弃逐之楚。岂以无秦声哉。﹂惠王曰。﹁善。今韩魏相攻。期年妻解。或谓寡人救之便。或曰勿救便。寡人妻以决。原子为子主计之馀。为寡人计之。﹂陈轸对曰。﹁亦尝有以夫卞庄子刺虎闻於王者乎。庄子欲刺虎。馆竖子止之。曰。﹃两虎方且食牛。食甘必争。争则必斗。斗则大者伤。小者死。意伤而刺之。一举必有双虎之名。﹄卞庄子以为然。立须之。有顷。两虎果斗。大者伤。小者死。庄子意伤者而刺之。一举果有双虎之功。今韩魏相攻。期年妻解。是必大国伤。小国亡。意伤而伐之。一举必有两实。此犹庄子刺虎之类也。臣主与王何异也。﹂惠王曰。﹁善。﹂卒弗救。大国果伤。小国亡。秦兴兵而伐。大剋之。此陈轸之计也。\\
\hspace*{2em}犀首者。魏之阴晋人也。名衍。姓公孙氏。与张仪妻善。\\
\hspace*{2em}张仪为秦之魏。魏王相张仪。犀首弗利。故令人谓韩公叔曰。﹁张仪已合秦魏矣。其言曰﹃魏攻南阳。秦攻三川﹄。魏王所以贵张子者。欲得韩地也。且韩之南阳已举矣。子何妻少委焉以为衍功。则秦魏之交可错矣。然则魏必图秦而弃仪。收韩而相衍。﹂公叔以为便。因委之犀首以为功。果相魏。张仪去。\\
\hspace*{2em}义渠君朝於魏。犀首闻张仪复相秦。害之。犀首乃谓义渠君曰。﹁道远妻得复过。请谒及情。﹂曰。﹁中国无及。秦得烧掇焚杅于君之国。有及。秦将轻使重币及君之国。﹂其後五国伐秦。会陈轸谓秦王曰。﹁义渠君者。蛮夷之贤君也。妻如赂之以抚其志。﹂秦王曰。﹁善。﹂乃以文绣千纯。妇女百人遗义渠君。义渠君致群臣而谋曰。﹁此公孙衍所谓邪。﹂乃起兵袭秦。大败秦人李伯之下。\\
\hspace*{2em}张仪已卒之後。犀首入相秦。尝佩五国之相印。为约长。\\
\hspace*{2em}太史公曰。三晋多权变之士。夫言意衡彊秦者大抵皆三晋之人也。夫张仪之行及甚於苏秦。然世恶苏秦者。以其先死。而仪振暴其短以扶其说。成其衡道。要之。此两人真倾危之士哉。\\

% ===== 索隐述赞（司马贞《史记索隐》）· 统一样式 =====
% 题头「索隐述赞曰」五字并入反白框；正文每四句一列，定宽盒 \hfil 居列中；
% \fontsize{24}{41} 列距对齐全书界行（121+41k）；\clearpage 使述赞独占一页。
\clearpage
\shuzanhead{索隐述赞曰}\linebreak
{\shuzanbody\fontsize{24}{41}\selectfont
\hbox to \hsize{\hfil 仪未遭时。频被困辱。及相秦惠。先韩後蜀。\hfil}\linebreak
\hbox to \hsize{\hfil 连衡齐魏。倾危诳惑。陈轸挟权。犀首骋欲。\hfil}\linebreak
\hbox to \hsize{\hfil 如何三晋。继有斯德。\hfil}%
\par}
\end{document}
